\documentclass[11pt]{article}
\usepackage[margin=1in]{geometry}
\usepackage{amsmath, amssymb, amsthm}
\usepackage{booktabs}
\usepackage{multirow}

\newcommand{\Usingle}{\mathcal{U}^{\text{single}}}
\newcommand{\UsNO}{\mathcal{U}^{s}_{NO}}
\newcommand{\UsO}{\mathcal{U}^{s}_{O}}
\newcommand{\UD}{\mathcal{U}^{D}}
\newcommand{\UDA}{\mathcal{U}^{D}_{A}}
\newcommand{\Umax}{\mathcal{U}^{\max}}
\newcommand{\Ugreedy}{\mathcal{U}^{\text{greedy}}_{A}}

\title{Computational Results: Dynamic Augmented Pooled Testing}
\author{Vladimir, Hector, Francisco}
\date{March 2026}

\begin{document}
\maketitle

\section{Strategy Definitions}

We consider six testing strategies for a population instance $J = (p_1, \ldots, p_n, u_1, \ldots, u_n)$
with budget $B$ and maximum pool size $G$:

\begin{enumerate}
    \item $\Usingle(J, B, G)$: Optimal strategy testing single individuals.
    Each test contains exactly one individual. Select the $\min(B, n)$
    individuals with highest $u_i \cdot q_i$.

    \item $\UsNO(J, B, G)$: Optimal \textbf{static non-overlapping} pooled tests
    (classical binary $+/-$ results). Choose $B$ disjoint pools of size $\leq G$
    before seeing any results. Individual $i$ is cleared if their pool tests negative.

    \item $\UsO(J, B, G)$: Optimal \textbf{static overlapping} pooled tests
    (classical binary $+/-$ results). Choose $B$ pools of size $\leq G$
    (pools may share individuals). Individual $i$ is cleared if \emph{any}
    pool containing $i$ tests negative.

    \item $\UD(J, B, G)$: Optimal \textbf{dynamic} pooled tests with
    \textbf{classical} binary results ($+/-$). The strategy adapts: the pool
    at step $k$ depends on previous outcomes. Test result is
    $r \in \{-, +\}$ where $-$ means ``no infected in pool.''

    \item $\UDA(J, B, G)$: Optimal \textbf{dynamic augmented} pooled tests.
    Same as $\UD$ but the test returns the \emph{exact count}
    $r = |t \cap Z| \in \{0, 1, \ldots, |t|\}$ of infected individuals in the pool.
    Strictly more informative than classical binary results.

    \item $\Umax(J, B, G) = \sum_{i=1}^{n} u_i \cdot q_i$: Upper bound
    achieved only with infinite budget.
\end{enumerate}

\begin{theorem}[Inequality Chain]
For any population instance $J$, budget $B \geq 1$, and pool size $G \geq 1$:
\[
    \Usingle \;\leq\; \UsNO \;\leq\; \UsO \;\leq\; \UD \;\leq\; \UDA \;\leq\; \Umax
\]
\end{theorem}

Additionally, we define a \textbf{myopic augmented greedy} strategy $\Ugreedy$
that at each step selects the pool maximizing immediate expected utility gain
$\prod_{i \in t} q_i \cdot \sum_{i \in t} u_i$, then updates beliefs via
Bayesian posterior inference using the augmented test count. By construction,
$\Usingle \leq \Ugreedy \leq \UDA$.


\section{Bayesian Posterior Update}

A key component of the augmented greedy is the Bayesian update after observing
$r = |t \cap Z|$ for a tested pool $t$. For each individual $i \in t$, the
posterior infection probability is:
\[
    p'_i = \mathbb{P}(Z_i = 1 \mid r) = \frac{\mathbb{P}(r \mid Z_i = 1) \cdot p_i}
    {\mathbb{P}(r \mid Z_i = 1) \cdot p_i + \mathbb{P}(r \mid Z_i = 0) \cdot q_i}
\]
where, letting $S = t \setminus \{i\}$:
\begin{align*}
    \mathbb{P}(r \mid Z_i = 1) &= \mathbb{P}\!\left(\sum_{j \in S} Z_j = r - 1\right)
    \quad \text{(Poisson-Binomial over $S$)} \\
    \mathbb{P}(r \mid Z_i = 0) &= \mathbb{P}\!\left(\sum_{j \in S} Z_j = r\right)
\end{align*}
For $i \notin t$, the test provides no information: $p'_i = p_i$.

\noindent Special cases:
\begin{itemize}
    \item $r = 0$: All pool members are proven healthy ($p'_i = 0$ for all $i \in t$).
    \item $r = |t|$: All pool members are proven infected ($p'_i = 1$ for all $i \in t$).
\end{itemize}


\section{Computational Results}

All values computed via brute-force exact enumeration over all $2^n$ infection profiles.

\subsection{Full Inequality Chain}

\begin{table}[h]
\centering
\caption{Comparison of all strategies across four instances.}
\label{tab:inequality}
\begin{tabular}{@{}lcccccc@{}}
\toprule
 & $\Usingle$ & $\UsNO$ & $\UsO$ & $\UD$ & $\UDA$ & $\Umax$ \\
\midrule
Instance 1 & 11.84 & 18.81 & 19.01 & 19.01 & 19.07 & 22.19 \\
Instance 2 & 8.00  & 8.60  & 8.60  & 8.67  & 8.91  & 12.40 \\
Instance 3 & 7.30  & 8.56  & 8.91  & 8.91  & 8.98  & 9.70 \\
Instance 4 & 1.70  & 3.29  & 3.29  & 3.29  & 3.29  & 4.25 \\
\bottomrule
\end{tabular}
\vspace{0.5em}

\footnotesize
Instance 1: $n{=}5$, $B{=}2$, $G{=}3$, $p{=}(0.05, 0.10, 0.15, 0.20, 0.08)$, $u{=}(4, 6, 3, 5, 7)$. \\
Instance 2: $n{=}4$, $B{=}2$, $G{=}2$, $p{=}(0.30, 0.40, 0.35, 0.25)$, $u{=}(5, 3, 4, 6)$. \\
Instance 3: $n{=}3$, $B{=}2$, $G{=}2$, $p{=}(0.10, 0.20, 0.30)$, $u{=}(5, 3, 4)$. \\
Instance 4: $n{=}5$, $B{=}2$, $G{=}3$, $p{=}(0.15, \ldots, 0.15)$, $u{=}(1, \ldots, 1)$.
\end{table}

\subsection{Augmented Benefit}

The augmented benefit measures how much the exact count information helps
compared to classical binary test results in the dynamic setting:
\[
    \text{Augmented Benefit} = \frac{\UDA - \UD}{\UD} \times 100\%
\]

\begin{table}[h]
\centering
\caption{Augmented benefit over classical dynamic testing.}
\label{tab:augmented-benefit}
\begin{tabular}{@{}lcccl@{}}
\toprule
 & $\UD$ & $\UDA$ & Benefit & Remark \\
\midrule
Instance 1 & 19.01 & 19.07 & $+0.28\%$ & Low infection rates \\
Instance 2 & 8.67  & 8.91  & $+2.74\%$ & High infection rates \\
Instance 3 & 8.91  & 8.98  & $+0.81\%$ & Mixed rates \\
Instance 4 & 3.29  & 3.29  & $+0.00\%$ & Uniform population \\
\bottomrule
\end{tabular}
\end{table}

Key observation: the augmented benefit is \textbf{largest when infection rates are high}
(Instance~2: $+2.74\%$). This aligns with the paper's intuition that augmented
information tests are more informative precisely when infection rates increase,
unlike classical protocols like Dorfman's which degrade.

For uniform populations (Instance~4), there is no benefit: symmetry means the
count provides no additional information about \emph{who} is infected.


\subsection{Greedy Performance}

\begin{table}[h]
\centering
\caption{Myopic augmented greedy vs.\ optimal augmented dynamic.}
\label{tab:greedy}
\begin{tabular}{@{}lcccl@{}}
\toprule
 & $\Ugreedy$ & $\UDA$ & Gap & \% of optimal gap captured \\
\midrule
Instance 1 & 19.053 & 19.066 & 0.013 & 99.8\% \\
Instance 2 & 8.876  & 8.909  & 0.033 & 96.3\% \\
Instance 3 & 8.966  & 8.982  & 0.016 & 99.1\% \\
Instance 4 & 3.287  & 3.287  & 0.000 & 100.0\% \\
\bottomrule
\end{tabular}
\vspace{0.5em}

\footnotesize
``\% of optimal gap captured'' $= (\Ugreedy - \Usingle) / (\UDA - \Usingle) \times 100\%$.
\end{table}

The myopic greedy with augmented Bayesian updates captures over 96\% of the
gap between individual testing and optimal in all instances. This suggests
that simple greedy policies with augmented posterior updates may be
near-optimal in practice.


\subsection{Simulation Example}

Consider Instance~1 with infection profile $Z = \{3\}$ (only individual~3
is infected).

\paragraph{Optimal strategy ($\UDA$).}
\begin{enumerate}
    \item Test pool $\{2,3,4\}$, result $r = 1$.
    \item Bayesian posteriors: $p'_2 = 0.34$, $p'_3 = 0.49$, $p'_4 = 0.17$.
    \item Test pool $\{0,1,4\}$, result $r = 0 \to$ \textsc{cleared}.
    \item Cleared: $\{0, 1, 4\}$, utility $= 17.0$.
\end{enumerate}

\paragraph{Myopic greedy.}
\begin{enumerate}
    \item Test pool $\{0,1,4\}$ (highest immediate score), result $r = 0 \to$ \textsc{cleared}.
    \item Test pool $\{2,3\}$, result $r = 1$.
    \item Cleared: $\{0, 1, 4\}$, utility $= 17.0$.
\end{enumerate}

Both achieve the same utility on this profile, but the optimal strategy
\emph{tests the riskier pool first} to gain information, while the greedy
\emph{tests the safest pool first} for immediate gain. The optimal strategy's
advantage appears in expectation over all possible infection profiles.


\end{document}
